% Transcription — fonds Alexandre Grothendieck, shelfmark 145, batch 1 (pages 1-20).
% Established from the facsimile archives/batches/145/batch-01.pdf (a mirror of
% https://grothendieck.umontpellier.fr/145.pdf), per the skill
% transcribe-grothendieck. The folder is thirty-seven pages; this is the first
% of two batches (batch 2, pages 21-37, is done separately).
%
% All twenty pages carry his hand and all twenty are transcribed; none is
% skipped and none is summarised as another's text. Page 3 is struck through
% whole by two long diagonals and is transcribed under that note. No title of
% his on pages 1-3; the section headings of pages 1 and 14 are the edition's
% and say so. Page 4 opens with his own title, under « [C] » and a circled 1.
% His own pagination: « 1 » on page 4, « 2 » on page 6, « 3 » on page 8,
% « 7 » on page 16, « 8 » on page 19 — recorded as \note{p. N de l'auteur};
% no number was seen on the pages between.
%
% What the batch is about. Pages 1-3 are three loose sheets on three kinds of
% paper, not in sequence: page 1, a group H of « u in G » for which there are
% finite-index pi', pi'' in pi and an iso v: D/pi' ~ D/pi'' covered by u on the
% universal cover D = U(C)~, its subgroups H_n; page 2, the automorphisms « à
% lacets » of Pi^_{0,3} commuting with sigma_0 and outer-commuting with rho,
% given by (g_1, p) subject to two boxed relations (1), (2) and the rewritten
% (1'); page 3 (struck), a diagram of classifying spaces B_{pi,Gamma}, B_pi,
% B_Gamma for phi: pi -> pi', automorphisms u' compatible with phi, and two
% short exact sequences of doubly-completed groups. Pages 4-14 (his pp. 1-3
% and following) are « Automorphismes extérieurs des groupes à lacets
% profinis Pi^ = Pi^_{0,3} »: presentation by l_0, l_1, l_inf with
% l_0 l_1 l_inf = 1; the circular automorphism rho and the transpositions
% sigma_i, rho sigma_i rho^-1 = sigma_rho(i); in Aut(pi) sigma_0 sigma_1 =
% int(l_1^-1) rho, so S_3 acts only as outer automorphisms (no fixed point of
% S_3 on P^1 - {0,1,inf}; Z/3 fixes ±j); complex conjugation tau with
% multiplier -1, giving an outer action of S_3 x {±1} = Gamma_{0,3}; outer
% automorphisms commuting with rho, normalised by u(l_0) = l_0^p and u rho =
% int(g) rho u, the boxed condition l_0^p . int(g) l_1^p . int(g rho(g))
% l_inf^p = 1; commutation with sigma_inf giving g l_1^q sigma_inf(g) = l_0^r
% with (q, r) invariant; tau gives g = l_1, q = -1, r = 1; the injection
% (gamma, chi): Pi -> pi/Z^ x Z^*, with Pi := Autext_lac(pi^_{0,3}, S_3) and
% phi(alpha).g = l_0^alpha g l_1^-alpha; then « quasi-square roots » epsilon_i
% of l_i, found after struck trials as eps_inf = l_0 sigma_inf l_inf etc.,
% sigma_i expressed through the eps_i, rho = eps_1^2 eps_0 eps_1 eps_0^-2 and
% cyclically, and the « inquiétant » (eps_0 eps_1 eps_inf)^2 = 1, then
% eps_0 eps_1 eps_inf = 1 and sigma_inf sigma_1 sigma_0 = 1 (this stretch is
% struck). Pages 14-20: « Critique des notations et de l'approche » — redo it
% in the discrete group with the fundamental groupoid of P^1(C) - {0,1,inf}:
% the Riemann sphere oriented so rho is rotation by +2pi/3 about the axis
% P, Pbar, sigma_i the symmetries about the axes (0,2), (1,-1), (inf,1/2), the
% margin's Möbius formulas; the extension E = pi_1(U_{0,3}, S_3 x Z/2; P)
% « à la Malgoire » with its product law; lifts sigma_i-dot, loops with
% l_inf l_1 l_0 = 1 (not l_0 l_1 l_inf); relations (2) sigma_0 sigma_1 =
% rho l_0 etc.; a presentation (4)-(6) by sigma_0, sigma_1, sigma_inf, rho,
% relation (7) for int(sigma_i), int(rho); eps_i redone (8) with rho eps_i
% rho^-1 = eps_rho(i), eps_0^2 = l_0 checked, a presentation of E^+ and the
% products (9); finally sigma~_i with (13), the inverse image Sigma~ of
% phi(S_3) as a semi-direct product of S_3 and Pi, the conjugation table
% (14), and tau_i = sigma_i sigma~_i with tau_i^2 = 1 (15)-(16), where the
% batch ends.
%
% Reading hazards fixed once here. S_3 is his round German S, set
% \mathfrak{S}_3; Pi^ and pi^ are both used on pages 2-14 for the profinite
% group, and Pi for the discrete one from page 15 on, as he writes them.
% « int(g) » is kept; « Aut ext » is set \mathrm{Autext}. His underlining is
% set in bold. Interlinear insertions of his are given as \marginal{} in
% place. Oblique margin notes are rendered where they read and otherwise as
% fragments with \ill.
%
% Readings to check first: page 1 almost throughout (the NB paragraphs and
% the H/pi sentence); the double-hatted script letters on page 3 (E, S, C —
% not distinguished); page 5 « (et même deux !) point fixe »; page 10 the
% second line of the tau system (« au moins, beta = -1 ») and « fonctions … beta »;
% page 11 the phrase on « quasi-… de ramification »; page 12 the exponent
% int(l_inf^alpha) and « doit écrire eps_0^2 = l_inf »; page 13 the margin
% « ces relations se déduisent … »; page 14 « Sans doute eps_0 eps_1 eps_inf =
% sigma_1 » (the right side is a short stroke); page 18 the exponents in the
% first system (8), which are overwritten; page 19 « relations des listes »
% and the formula introducing (7); page 20 « qui transforme les lacets … » and
% the three margin notes; every oblique margin note on pages 2, 7, 8, 9, 10,
% 17, 19, 20.
%
% Pass: Opus 5.5 (claude-opus-5-5), 2026-09-24 — first pass, unchecked against the pages by a human.
\documentclass[11pt,a4paper]{article}
% Le préambule commun à toutes les transcriptions du fonds Grothendieck.
%
% Il tient en une idée : **l'incertitude se compose, elle ne se résout pas.**
% Une transcription qui lisse un mot douteux a détruit la seule chose pour
% laquelle on la faisait — un lecteur doit pouvoir savoir, à chaque endroit,
% ce qui était sur le feuillet et ce qui a été ajouté. Les six macros qui
% suivent sont donc l'appareil critique, et non de la décoration.
%
% Les mêmes commandes sont reconnues par `scripts/render.mjs`, qui produit la
% vue de lecture du site. Une macro ajoutée ici doit l'être aussi là-bas,
% faute de quoi le rendu échouera bruyamment — ce qui est voulu : un rendu
% silencieusement faux est pire qu'un rendu absent.

\ProvidesPackage{grothendieck}[2026/08/08 Transcriptions du fonds Grothendieck]

\usepackage{amsmath,amssymb,amsthm}
\usepackage{mathrsfs}
\usepackage{stmaryrd}
% Les diagrammes commutatifs sont partout dans ce fonds : c'est la langue
% même de la géométrie algébrique de Grothendieck, et une transcription qui
% les rendrait en prose ne transcrirait rien.
\usepackage{tikz-cd}
% Français par défaut — c'est la langue des feuillets — mais l'anglais est
% chargé aussi : la traduction et le résumé sont des documents anglais, et la
% typographie française (espaces avant les deux-points) y serait une faute.
% Ces documents basculent par \selectlanguage{english} en tête de corps.
\usepackage[english,french]{babel}
\usepackage{fontspec}
\usepackage{geometry}
\geometry{a4paper,margin=25mm,marginparwidth=28mm}
\usepackage{marginnote}
\usepackage{xcolor}
% ulem, not soul. `soul`'s \st and \ul are built on a character-by-character
% scan that XeLaTeX's font machinery breaks: the document fails with an
% undefined control sequence rather than degrading. ulem does the same two jobs
% and is Unicode-engine safe. `normalem` stops it hijacking \emph, which the
% transcriptions use for Grothendieck's own emphasis.
\usepackage[normalem]{ulem}
\usepackage{hyperref}
\hypersetup{colorlinks=true,linkcolor=black,urlcolor=black,pdfborder={0 0 0}}

% --- Métadonnées du lot -----------------------------------------------------
% Reprises telles quelles par le script de rendu, qui les affiche en tête de
% la vue de lecture. Elles fixent ce que le fichier prétend couvrir : sans
% elles, une transcription flotte sans point d'attache dans un fonds de seize
% mille feuillets.

\newcommand{\folder}[1]{\gdef\grothFolder{#1}}
\newcommand{\batch}[1]{\gdef\grothBatch{#1}}
\newcommand{\leaves}[2]{\gdef\grothFirst{#1}\gdef\grothLast{#2}}
\newcommand{\foldertitle}[1]{\gdef\grothFoldertitle{#1}}
\gdef\grothFolder{?}\gdef\grothBatch{?}\gdef\grothFirst{?}\gdef\grothLast{?}\gdef\grothFoldertitle{}

% --- Le résumé --------------------------------------------------------------
%
% Il ouvre la lecture modernisée, sous le titre « Résumé », et il est composé
% en petit. Ce n'est pas un ornement : c'est le seul passage du document qui
% ne soit pas des mathématiques, et un lecteur doit pouvoir le reconnaître
% avant de l'avoir lu — puis le sauter s'il connaît déjà le sujet, ou s'y
% arrêter s'il ne le connaît pas. Le filet qui le suit marque la reprise du
% corps.

\newenvironment{resume}
  {\small\setlength{\parskip}{0.45em}}
  {\par\medskip\hrule\medskip}

% --- Mots-clés ---------------------------------------------------------------
%
% En anglais, en clôture du résumé de la lecture modernisée : le vocabulaire
% moderne sous lequel chercher ce que les feuillets construisent. Le manifeste
% du site (`npm run manifest`) les extrait pour étiqueter la cote entière —
% cette ligne est la source unique des tags, il n'y a pas de fichier à côté.
\newcommand{\keywords}[1]{\par\smallskip\noindent
  {\footnotesize\textsc{Keywords} --- \textit{#1}}\par}

% --- L'appareil critique ----------------------------------------------------

% Le numéro de feuillet, en marge. C'est l'ancre qui permet de régler un
% désaccord sur une formule en regardant *un* feuillet plutôt qu'une chemise
% de six cents. Le site s'en sert aussi pour tourner les pages du fac-similé
% quand on fait défiler la transcription.
\newcommand{\leaf}[1]{%
  \marginnote{\footnotesize\color{gray!70}\textsf{#1}}%
  \label{leaf:#1}\ignorespaces}

% Illisible : on ne devine pas. Un mot inventé qui se lit comme les autres est
% le pire résultat possible.
\newcommand{\ill}{\textcolor{red!60!black}{[\,\ldots\,]}}

% Lecture douteuse : le mot est donné, et signalé comme incertain.
\newcommand{\uncertain}[1]{\uline{#1}}

% Ajout de l'éditeur — une lettre manquante, un mot sous-entendu.
\newcommand{\add}[1]{\textcolor{blue!50!black}{[#1]}}

% Biffé par Grothendieck lui-même. Ce qu'il a rayé fait partie du document :
% une rature dit souvent où la pensée a changé de direction.
\newcommand{\struck}[1]{\sout{#1}}

% Note du transcripteur, sur l'état matériel du feuillet.
\newcommand{\note}[1]{\par\smallskip{\small\color{gray!60!black}\textit{[#1]}}\par\smallskip}

% Note marginale de Grothendieck, distincte du corps du texte.
\newcommand{\marginal}[1]{\par\smallskip\noindent\hspace{1em}%
  {\small\color{gray!60!black}#1}\par\smallskip}

% --- Titre ------------------------------------------------------------------

\AtBeginDocument{%
  \begin{center}
    {\large\bfseries Fonds Alexandre Grothendieck --- cote n\textsuperscript{o}~\grothFolder}\\[2pt]
    {\itshape\grothFoldertitle}\\[4pt]
    {\small feuillets \grothFirst--\grothLast\ (lot \grothBatch)}\\[2pt]
    {\footnotesize\color{gray!60!black}Transcription établie d'après le fac-similé numérisé
     par l'Université de Montpellier.\\
     Les lectures incertaines sont soulignées, les illisibles notées [\,\ldots\,],
     les ajouts entre crochets.}
  \end{center}
  \vspace{1em}}

\endinput


\folder{145}
\batch{1}
\pages{1}{20}
\dating{1981}
\watermark{Édition de démonstration}
\foldertitle{Printemps 81 (vieille rédaction) : notes manuscrites (s.d.)}

\begin{document}

\section{Feuillets épars : le groupe $H$, automorphismes à lacets, et un diagramme de $B_{\pi,\Gamma}$}

\note{ce titre est de l'édition. Les pages 1 à 3 sont trois feuillets de papiers
différents (papier blanc, papier listing perforé, papier à musique), sans
pagination de sa main ; ils ne se suivent pas}

\page{1}

\struck{bien déterminé \ill{} $\overline{D}_{J,\alpha}$ de $\overline{D}_J$, cf. \ill{}}

\begin{tikzcd}
\overline{U}_{J',\alpha'} \arrow[r] \arrow[d] & \overline{U}_{J,\alpha} \arrow[r] \arrow[d] & \overline{U} \arrow[d] \\
U_{J',\alpha'} \arrow[r] & U_{J,\alpha} \arrow[r] & U
\end{tikzcd}

\note{la ligne du haut et le diagramme sont barrés de grandes croix}

\begin{tikzcd}
D = U(\mathbb{C})^{\sim} \arrow[rr] \arrow[dr, "\pi"'] & & \widetilde{U}(\mathbb{C}) \arrow[dl, "\hat\pi"] \\
 & U(\mathbb{C}) &
\end{tikzcd}

\[
\pi \subset G
\]
\[
H = \Bigl\{ u \in G \Bigm| \exists\; \underbrace{\pi' \subset \pi}_{\substack{\text{ss-gr.} \\ \text{d'indice fini}}},\;
\underbrace{\pi'' \subset \pi}_{\text{id.}},
\]
iso $v : D/\pi' \simeq D/\pi''$, avec

\begin{tikzcd}
D \arrow[r, "u"] \arrow[d, "p'"'] & D \arrow[d, "p''"] \\
D/\pi' \arrow[r, "\sim"'] & D/\pi''
\end{tikzcd}

commutatif

\note{après $\widetilde{U}(\mathbb{C})$, un signe biffé. L'accolade fermante de $H$ n'est pas écrite ; un signe « $\sim$ » est
porté au centre du carré}

\textbf{NB} On aura, pour $v$, $\pi'$, $\pi''$, $u$ comme dessus, un iso
\[
v_* : \pi' \longrightarrow \pi''
\]
(induit par $u$ \uncertain{on aura d'ailleurs} $\pi'$ et $\pi''$ de même indice
dans $\pi$.
\marginal{\uncertain{$u$ est déterminé par} $v_*$}

\textbf{NB} $v$ est connu quand on connaît $v_*$ (cas \uncertain{connexe} \ill{}),
et quand $v$ \uncertain{est} \uncertain{connu}, on ne \uncertain{connaît} \ill{}
\uncertain{$u$} qu'à \ill{} multiplication : \uncertain{à gauche} par \uncertain{un}
\uncertain{élt} de $\pi''$.
\[
\begin{array}{c}
\pi \\
\cup \quad \cup \\
\pi' \simeq \pi'' \\
\pi'_i \simeq \pi''_i
\end{array}
\]

$H$ est un groupe $\supset \pi$, et $H/\pi$ (\uncertain{classes} à droite)
s'identifie à un sous-\ill{} \uncertain{de} l'un des « groupoïdes »
d'isom. \struck{\ill{}} \uncertain{H-groupoïde} d'indice \ill{} de $\pi$ de
\ill{} indice
\[
p'' u = v p', \qquad \struck{p'' u \gamma =}, \qquad
p''(\gamma u) = (p''\gamma) u = p'' u = v p'
\]
\[
\struck{\pi \subset H}
\]
$H_n$ = sous-ens. de $H$ formé des $u$ pour lesquels on peut prendre
$\pi'$, $\pi''$ d'indice $n$.

$H_n$ est un ss-groupe, et $\pi$ d'indice

\note{la page s'arrête sur ces mots}

\page{2}

Les autom. \add{à lacets} de $\hat\Pi_{0,3} = \hat\pi$ commutant \struck{\ill{}} à
$\sigma_0$, et commutant extérieurement à $\rho$, sont donnés par
\[
\begin{aligned}
u(l_1) &= \mathrm{int}(g_1)\, l_1^{p} \\
u(l_\infty) &= \mathrm{int}(\underbrace{\sigma_0(g_1)}_{g_\infty})\, l_\infty^{p}
\end{aligned}
\]
avec
\[
g_1 \in \pi \bmod L_1, \qquad p \in \hat{\mathbb{Z}}^{*} \quad \text{satisfaisant}
\]
la condition (où $\beta, \gamma \in \hat{\mathbb{Z}}$ sont uniquement déterminés
par $g_1, p$)
\[
(1) \qquad \boxed{l_\infty^{p}\,.\,\mathrm{int}(h_1)(l_1^{p})\,.\,
\mathrm{int}\bigl(h_1\, l_1^{-\gamma}\, \rho^{-1}(h_1)\bigr)(l_0^{p}) = 1}
\]
\[
\text{et} \quad (2) \qquad \boxed{\rho(h_1)\bigl(l_\infty^{-\gamma}\, h_1\, l_1^{-\gamma}\bigr)\rho^{-1}(h_1) = l_0^{\beta}}
\]
où
\[
h_1 = \sigma_0(g_1)^{-1}\, g_1
\]

\note{l'indice de $u(l_1)$ est surchargé ; « à lacets » est un ajout
interlinéaire}

\marginal{NB On aura $u(l_0) = \mathrm{int}\bigl(g_1\, l_1\, \rho^{-1}(h_1)\bigr).\, l_0^{p}$,
avec $g_0$ \uncertain{sous} le premier facteur}

\marginal{\ill{} à \ill{} près ! \ill{}}

\marginal{NB. On doit \ill{} \ill{} \uncertain{$\beta$} \ill{} $\Pi_{0,3}$ \ill{}}

\note{la première note marginale est écrite en oblique dans l'angle supérieur
droit ; les deux autres, dans la marge gauche, ne se lisent que par
fragments}

La deuxième relation s'écrit aussi
\[
h_1\, l_1^{-\gamma}\, \rho^{-1}(h_1) = \struck{l_\infty} l_\infty^{\gamma}\, \rho(h_1)^{-1}\, l_0^{\beta}
\]
et en portant dans la première relation celle-ci prend la forme
\[
(1') \qquad \boxed{l_\infty^{p}\,.\,\mathrm{int}(h_1)(l_1^{p})\,.\,
\mathrm{int}\bigl(l_\infty^{\gamma}\, \rho(h_1)^{-1}\bigr)(l_0^{p}) = 1}
\]

\textbf{NB} Si on change $g_1$ en $g'_1 = g_1\, l_1^{\mu}$, $h_1$ est remplacé
par
\[
h'_1 = l_\infty^{-\mu}\, h_1\, l_1^{\mu}
\]
et la relation (1') donne, \uncertain{avec} $\mathrm{int}(l_\infty^{-\mu})$ :
\[
l_\infty^{p}\,.\,\mathrm{int}(\uncertain{h'_1})(l_1^{p})\,.\,
\mathrm{int}\bigl(\struck{\ill{}}\; l_\infty^{\gamma - \mu}\, \rho^{-1}(h_1'^{-1})\bigr)(l_0^{p}) = 1
\]
\uncertain{(i.e. $\gamma' = \gamma \ill{}$)} et on a \uncertain{aussi} la relation
correspondante (2), avec $\beta' = \beta$.

\note{dans la dernière formule, les exposants de $l_\infty$ et la place de
$\rho^{-1}$ sont surchargés ; la parenthèse « (i.e. $\gamma' = \gamma \ldots$ »
se termine par un signe gratté}

\page{3}

\begin{tikzcd}
\pi \arrow[d, "\varphi"'] & \Gamma & B_{\pi,\Gamma} \arrow[d] & B_{\pi} \arrow[l] \arrow[d] & \widetilde{B}_{\pi} \arrow[l] \\
\pi' & \Gamma & B_{\pi',\Gamma} \arrow[d, no head] & B_{\pi'} \arrow[l] \arrow[d] & \\
 & & B_{\Gamma} & \widetilde{B}_{\Gamma} \arrow[l] &
\end{tikzcd}

\begin{tikzcd}
B_{\pi,\Gamma} \arrow[r, no head] \arrow[d, no head] & B_{\pi,\Gamma_i} \arrow[d, no head] & B_{\pi} \arrow[l] \arrow[d, no head] \\
B_{\pi',\Gamma} \arrow[r, no head] \arrow[d, no head] & B_{\pi',\Gamma_i} \arrow[d, no head] & B_{\pi'} \arrow[l] \arrow[d, no head] \\
B_{\Gamma} \arrow[r, no head] & B_{\Gamma_i} \arrow[r, no head] & \widetilde{B}_{\Gamma}
\end{tikzcd}

\note{le second diagramme est écrit à droite du premier, en plus petit ; ses
traits sont sans flèche sauf les deux horizontales de droite des deux
premières lignes}

$u$ automorphisme de $\pi$, compatible avec $\varphi$

$\exists !$ $u'$ automorphisme de $\pi'$, compatible \uncertain{avec} $\varphi$,

tel que $u'\varphi(g) = \varphi(ug)$ \quad $\forall g \in \pi$
\[
\struck{u''\varphi} \qquad u'\varphi u^{-1} = \varphi
\]
\[
u'\varphi = \varphi u \qquad u'\varphi = v'
\]
\[
v'\varphi = \varphi u
\]
\[
u'\varphi u^{-1} = \mathrm{int}(g)\, \varphi u
\]
\[
u'\varphi = \struck{\varphi u}\; \mathrm{int}(g)\, \varphi(u)
\]

\note{à gauche des deux dernières lignes, un $u$ au-dessus d'un trait
vertical}

Soit \struck{$\mathcal{E}$}
\[
\mathcal{E}_{\varphi} = \bigl\{ u \in \mathrm{Aut}(\pi) \bigm| \exists\, u' \in \mathrm{Aut}(\pi')\; \ldots
\]

\begin{tikzcd}
\pi \arrow[d, "\varphi"'] & \hat{\hat{\mathcal{E}}}(\pi, I') \arrow[d, "\wr\; \Psi_{I'}"] & 1 \arrow[r] & \pi \arrow[r] \arrow[d, "f_{\pi}"] & \hat{\hat{\mathfrak{S}}}(\pi, I') \arrow[r] \arrow[d, "f_{\hat{\hat{\mathfrak{S}}}}"] & {} \\
\pi' & \hat{\hat{\mathfrak{S}}}(\pi') \arrow[d, "p"] & 1 \arrow[r] & \pi' \arrow[r] & \hat{\hat{\mathcal{E}}}(\pi') \arrow[r] & {} \\
 & \hat{\hat{\mathcal{C}}}(\pi') & & & &
\end{tikzcd}

\note{sous $\hat{\hat{\mathcal{E}}}(\pi')$, entre parenthèses :
« $\simeq \hat{\hat{\mathcal{C}}}(\pi, I')$ ». Une flèche courbe va de
$\hat{\hat{\mathcal{E}}}(\pi, I')$ à $\hat{\hat{\mathcal{C}}}(\pi')$, marquée
« compatible \uncertain{avec} $\varphi$ ». Les lettres rondes, surmontées de
deux chapeaux, sont mal distinguées l'une de l'autre ; les lectures
$\mathcal{E}$, $\mathfrak{S}$ et $\mathcal{C}$ sont douteuses}

\note{toute la page est barrée de deux longs traits obliques}

\section{Automorphismes extérieurs des groupes à lacets profinis $\hat\Pi = \hat\Pi_{0,3}$}

\note{ce titre est le sien ; il ouvre la page 4, que coiffent « [C] » et un
« ① » cerclé}

\page{4}
\note{p. 1 de l'auteur}

\marginal{et structure d'extension de $\mathcal{E} = \mathcal{E}_{0,3}$}

On prend la présentation par les générateurs $l_0, l_1, l_\infty$ liés par
\[
l_0\, l_1\, l_\infty = 1
\]
Un endomorphisme de $\Pi_{0,3}$ est donc donné par des éléments
$l'_0, l'_1, l'_\infty$ satisfaisant la même relation — sauf erreur ce sera
un automorphisme s'ils engendrent $\pi$ (topologiquement).

Soit $\rho$ l'automorphisme défini par la permutation circulaire
\[
\rho(l_0) = l_1, \quad \rho(l_1) = l_\infty, \quad \rho(l_\infty) = l_0
\]
(donc aussi $\rho^3 = \mathrm{id}$). Considérons aussi la transposition
\struck{On s'intéresse aux}
\[
\sigma_\infty(l_0) = l_1, \quad \sigma_\infty(l_1) = l_0, \quad
\sigma_\infty(l_\infty) = (l_1 l_0)^{-1} = l_0^{-1} l_1^{-1}
\quad \bigl(= \mathrm{int}(l_0^{-1}).\, l_\infty\bigr)
\]
on a évidemment $\sigma_\infty^2 = \mathrm{id}$. On définit de même
$\sigma_0$, $\sigma_1$ (échangeant $l_1$ et $l_\infty$ resp. $l_\infty$ et $l_0$),
de sorte que
\[
\rho\, \sigma_i\, \rho^{-1} = \sigma_{\rho(i)}
\]
en désignant par $\rho$ aussi la permutation circulaire sur $0, 1, \infty$. On
considère $\rho$ et les $\sigma_i$ comme des éléments du groupe
$\mathfrak{S}_3$ des permutations de l'ens. $\{0, 1, \infty\}$, qu'on se
propose de faire opérer sur $\pi$, en respectant la structure locale (avec
\marginal{\uncertain{correction}} multiplicateur 1). On a bien une
\marginal{vraie} représentation du ss-groupe
$\mathfrak{S}_3^{+} \simeq \mathbb{Z}/3\mathbb{Z}$ engendré


\page{5}

par $\rho$, mais pas de $\mathfrak{S}_3$, car on a p.ex
\[
\sigma_0\sigma_1 \;\bigl(= \sigma_1\sigma_\infty = \sigma_\infty\sigma_0\bigr) = \rho
\]
\[
\struck{[\sigma_0, \sigma_1] \ill{}} = \struck{(\sigma_0\sigma_1)^2 = \rho^2}
\qquad \text{dans } \mathfrak{S}_3
\]

\note{la seconde ligne est hachurée ; on y devine un commutateur
$[\sigma_0, \sigma_1]$ et $(\sigma_0\sigma_1)^2 = \rho^2$}

mais dans $\mathrm{Aut}(\pi)$, \struck{$\sigma_0\sigma_1$}
\[
\begin{aligned}
\sigma_0\sigma_1(l_0) &= \sigma_0(l_\infty) = l_1 = \rho(l_0) \\
\sigma_0\sigma_1(l_1) &= \sigma_0(l_\infty^{-1} l_0^{-1}) = \sigma_0(l_\infty)^{-1}\sigma_0(l_0)^{-1}
= l_1^{-1}\,(l_1^{-1} l_\infty^{-1})^{-1} = l_1^{-1}\, l_\infty\, l_1 \\
\sigma_0\sigma_1(l_\infty) &= (l_1^{-1} l_\infty l_1)^{-1}\, l_1^{-1}
= l_1^{-1} l_\infty^{-1} \struck{l_1} = l_1^{-1}\,(l_\infty^{-1} l_1^{-1})\, l_1 = l_1^{-1}\, l_0\, l_1
\end{aligned}
\]
donc
\[
\sigma_0\sigma_1 = \mathrm{int}(l_1^{-1})\,\rho \qquad \text{dans } \mathrm{Aut}(\pi)
\]
et de même
\[
\sigma_1\sigma_\infty = \mathrm{int}(l_\infty^{-1})\,\rho, \qquad
\sigma_\infty\sigma_0 = \mathrm{int}(l_0^{-1})\,\rho
\]

On a une représentation de $\mathfrak{S}_3$ dans le groupe
\textbf{extérieur} à lacets seulement. Notons qu'on ne peut espérer réaliser
cette opération extérieure par une vraie opération (en \uncertain{restant} pour
le groupe discret) vu que $\mathfrak{S}_3$ opérant sur
$\mathbb{P}^1_{\mathbb{C}} \setminus \{0, 1, \infty\}$ \uncertain{n'a pas}
de pt fixe. Le fait qu'on ait pu faire opérer
\marginal{\uncertain{sur $\Pi$ lui-même}} $\mathbb{Z}/3\mathbb{Z}$ \uncertain{provient}
du fait que son action sur $\mathbb{P}^1_{\mathbb{C}} - \{0, 1, \infty\}$ admet
un \uncertain{(et même} deux !) point fixe, qui \uncertain{peut} être
\struck{\ill{}} \uncertain{pris} comme origine pour calculer le $\pi_1$
\marginal{, savoir $\pm j$ (racine primitive troisième de 1)}.
Signalons aussi \uncertain{qu'on} \uncertain{écrira} dans le groupe
\marginal{$\Gamma_{0,3}$ opérant sur $\Pi_{0,3}$} de Teichmüller) la
« conjugaison complexe »

\note{les ajouts « sur $\Pi$ lui-même », « savoir $\pm j$ (racine primitive
troisième de 1) » et « $\Gamma_{0,3}$ opérant sur $\Pi_{0,3}$ » sont
interlinéaires, reliés au texte par des traits ; la parenthèse ouvrante
avant « de Teichmüller » manque}

\page{6}
\note{p. 2 de l'auteur}

$\tau$, dont l'opération sur $\Pi_{0,3}$ peut se décrire par
\[
\tau(l_0) = l_0^{-1}, \qquad \tau(l_1) = l_1^{-1}, \qquad
\tau(l_\infty) = l_1\, l_0 = \mathrm{int}(l_1)\, l_\infty^{-1}
\]
C'est donc un autom. du \uncertain{groupe} à lacets, de
\uncertain{multiplicateur} $-1$, satisfaisant
\[
\tau^2 = 1
\]

On vérifie qu'il commute à l'action de $\mathfrak{S}_3$ mod autom.
extérieurs — ce qui donne une op. extérieure de
$\mathfrak{S}_3 \times \{\pm 1\} = \Gamma_{0,3}$, qui \uncertain{n'est
autre} que l'action canonique.

On s'intéresse maintenant aux autom. extérieurs qui commutent à $\rho$.
Notons que tout autom. extérieur provient d'un autom $u$ qui normalise
$L_0$ ($= \hat{\mathbb{Z}}$ engendré par $l_0$)
\marginal{qui normalise la classe de conj. de $l_0$}
et $u$ est unique mod. conjugaison par un $\lambda \in L_0$. Si
$p \in \hat{\mathbb{Z}}^{*}$ \marginal{$= \chi(u)$,} \struck{est} \uncertain{est} le
multiplicateur de $u$, on aura donc
\[
(*) \qquad u(l_0) = l_0^{p}
\]
Soit $g \in \pi$ tel que
\[
(**) \qquad u\rho = \struck{\rho u}\; \mathrm{int}(g)\, \rho u ,
\]
ce qui détermine $g$ de façon unique, \uncertain{u} \uncertain{étant}
\ill{}

\note{« qui normalise la classe de conj. de $l_0$ » est écrit sous la ligne,
en petit, sans marque d'insertion ; « $= \chi(u)$, » est ajouté au-dessus
de « $p \in \hat{\mathbb{Z}}^{*}$ »}

\page{7}

choisi — si on \textbf{multiplie} remplace $u$ par $\mathrm{int}(\lambda)u$
avec $\lambda = l_0^{\alpha} \in L_0$ $(\alpha \in \hat{\mathbb{Z}}^{*})$, alors $g$
est remplacé par
\[
\lambda\, g\, \rho(\lambda)^{-1} = \mathbf{l_0^{\alpha}\, g\, l_1^{-\alpha}}.
\]
Ceci posé,

les relations $(*)$ et $(**)$ déterminent $u$ en termes de
$p \in \hat{\mathbb{Z}}^{*}$ et de $g \in \pi$, par
\[
\begin{cases}
u(l_0) = l_0^{p} \\
u(l_1) = u\rho(l_0) = \mathrm{int}(g)\,\rho\, \underbrace{u(l_0)}_{l_0^{p}} = \mathrm{int}(g)\, l_1^{p} \\
u(l_\infty) = u\rho(l_1) = \mathrm{int}(g)\, \rho\, \underbrace{u(l_1)} = \mathrm{int}(g)\, \underbrace{\rho\, \mathrm{int}(g)\, \rho^{-1}}_{\mathrm{int}(\rho(g))}\, \underbrace{(\rho\, l_1^{p})}_{l_\infty^{p}} = \struck{\mathrm{int}(g\rho(g))\ill{}}
\end{cases}
\]
\marginal{$\mathrm{int}\, g\rho(g))\, l_\infty^{p}$}

\note{sous $u(l_1)$ dans la troisième ligne, l'accolade porte
$\mathrm{int}(g)\, l_1^{p}$ ; le résultat biffé en bout de ligne est récrit au-dessus,
à droite : $\mathrm{int}(g\rho(g))\, l_\infty^{p}$ — la parenthèse ouvrante manque}

\marginal{NB On sait que \ill{} \uncertain{automorphismes} \ill{} respectant la
\uncertain{structure} \ill{} à lacets \ill{} multiplicateur $p$}

\note{note oblique de la marge gauche, en regard du système}

Il faut vérifier, pour qu'on ait bien un endom. de $\pi$, la relation
$l'_0\, l'_1\, l'_\infty = 1$ i.e.
\[
\boxed{l_0^{p}\,.\,\mathrm{int}(g)\, l_1^{p}\,.\,\mathrm{int}\bigl(g\rho(g)\bigr)\, l_\infty^{p} = 1}
\]
ou encore
\[
l_0^{p}\,\bigl(g\, l_1^{p}\, g^{-1}\bigr)\bigl(g\rho(g)\, l_\infty^{p}\, \rho(g)^{-1} g^{-1}\bigr) = 1
\]
i.e.
\[
l_0^{p}\, g\, l_1^{p}\, \rho(g)\, l_\infty^{p}\, \rho(g)^{-1}\, g^{-1} = 1
\]
soit encore, en faisant passer le $g^{-1}$ à gauche
\[
\struck{l_0^{p}\,.\,\mathrm{int}(g)\bigl(l_1^{p}\,.\,\mathrm{int}(\rho(g))(l_\infty^{p})\bigr) = 1}
\]
\[
\mathrm{int}(g^{-1})\, l_0^{p}\,.\, l_1^{p}\,.\,\mathrm{int}(\rho(g))(l_\infty^{p}) = 1
\]

Les relations $(*)$ et $(**)$ sont vérifiées par construction (pour $(**)$,
il suffit \struck{\ill{}} de la vérifier pour la valeur aux deux arguments
$l_0, l_1$).

\note{« $(*)$ et $(**)$ » : le second signe est surchargé}

\page{8}
\note{p. 3 de l'auteur}

Mais comment vérifier que $(p, g)$ définissent non seulement un endomorphisme,
mais un automorphisme ? Ça semble délicat …

Dans le cas de $\tau$, on a $p = -1$, et $g = l_1$ :
\[
\tau\rho = \mathrm{int}(l_1)\, \rho\, \tau
\]

On veut exprimer la commutation de $u$
\marginal{respectant la structure à lacets,}
\struck{\ill{}} autom. ext.) non seulement à $\rho$ i.e.
$\mathfrak{S}_3^{+}$, mais à $\mathfrak{S}_3$ tout entier. Il suffit
d'exprimer la commutation mod autom. \uncertain{intérieurs} à $\sigma_\infty$,
i.e. d'écrire qu'il existe $h \in \pi$ (\uncertain{également} unique) tel que
\[
u\,\sigma_\infty = \mathrm{int}(h)\, \struck{\ill{}}\, \sigma_\infty\, u ,
\]
et il suffira de l'exprimer pour les \textbf{arguments} $l_0$ et $l_1$. On
trouve les conditions
\[
\underbrace{u\,\sigma_\infty(l_0)}_{\substack{\| \\ u(l_1) \\ \mathrm{int}(g)\, l_1^{p}}}
= \mathrm{int}(h)\, \underbrace{\sigma_\infty(l_0^{p})}_{l_1^{p}}
\qquad \text{i.e.} \qquad \mathrm{int}(g^{-1}h)\, l_1^{p} = l_1^{p}
\]
\[
\text{i.e.} \qquad g^{-1}h \in L_1 \quad \text{i.e.} \quad h = g\, l_1^{q},\; q \in \hat{\mathbb{Z}}
\]
et
\[
\underbrace{u\,\sigma_\infty(l_1)}_{l_0} = \mathrm{int}(h)\, \sigma_\infty\, \underbrace{u(l_1)}_{\mathrm{int}(g)\, l_1^{p}}
\qquad \text{i.e.}
\]
\[
l_0^{p} = \mathrm{int}\bigl(h\, \sigma_\infty(g)\bigr)\, l_0^{p} \qquad \text{soit}
\]
\[
l_0^{p} = \mathrm{int}\bigl(g\, l_1^{q}\, \sigma_\infty(g)\bigr)\, l_0^{p}
\]
ce qui équivaut à
\[
\struck{\mathrm{int}(g\, l_1^{q}\, \sigma_\infty g)\, L_1 = L_0}
\qquad \text{i.e. la condition}
\]
\[
\exists\, q \quad \text{tel que} \quad g\, l_1^{q}\, \sigma_\infty(g) \in L_0
\]

\note{sous $u\sigma_\infty(l_0)$, un signe $=$ vertical et $u(l_1)$, puis
l'accolade $\mathrm{int}(g)\, l_1^{p}$ ; le premier membre est surchargé.
Dans la seconde condition, l'accolade sous $u\sigma_\infty(l_1)$ porte $l_0$,
et le premier membre écrit est $l_0^{p}$}

\marginal{NB Ceci \uncertain{implique} qu'il \uncertain{existe} \ill{}
\ill{} $\Gamma$ \ill{} très \ill{} \ill{} sous-groupe extérieur \ill{}
\uncertain{qui} \ill{}}

\note{note oblique de la marge gauche, en regard de « $\mathfrak{S}_3^{+}$ »}

\page{9}

on aura
\[
\exists\, q, r \in \hat{\mathbb{Z}}^{*} \quad \text{tels que} \quad
\boxed{g\, l_1^{q}\, \sigma_\infty(g) = l_0^{r}}
\]
(condition qui ne dépend que de $g$, pas de $p$ ; si $g$ est remplacé par
$g' = l_0^{\alpha}\, g\, l_1^{-\alpha}$, on aura
\[
\sigma_\infty(g') = l_1^{\alpha}\, \sigma_\infty(g)\, l_0^{-\alpha},
\]
et la relation
\[
g'\, l_1^{q'}\, \sigma_\infty(g') = l_0^{r'}
\]
s'écrit
\[
l_0^{\alpha}\, g\, l_1^{-\alpha}\, l_1^{q'}\, l_1^{\alpha}\, \sigma_\infty(g)\, l_0^{-\alpha} = l_0^{r'}
\]
i.e. en simplifiant
\[
g\, l_1^{q'}\, \sigma_\infty(g) = l_0^{r'}
\]
donc $q = q'$, $r = r'$ i.e. $(q, r)$ ne dépend que de l'autom. ext. défini
par $u$, et pas du \struck{choix} \uncertain{des} représentants dans le
normalisateur de $L_0$.

\note{la parenthèse ouverte avant « condition » n'est pas refermée}

\marginal{il vaut mieux écrire la relation sous la forme
$\sigma_\infty(g^{-1}) = l_0^{\beta}\, \struck{\ill{}}\, g\, l_1^{\gamma}$
(conditions \uncertain{de} définition) \ill{} on trouve \ill{} \ill{}
(i.e. $\beta = r$, $\gamma = -q$) donc
$\sigma_\infty(g^{-1}) = l_0^{\beta}\, g\, l_1^{\gamma}$}

\note{la note est écrite en oblique, de bas en haut, dans la marge gauche,
en regard des premières lignes de la page}

Dans le cas de $\tau$, où $g = l_1$, la relation devient
\[
l_1\, l_1^{q}\, l_0 = l_0^{r} \qquad \text{donc} \qquad q = -1,\; r = 1 .
\]

On sait que, \uncertain{travaillant} avec les groupes discrets, en plus de
l'autom. ext. identique, il n'y a qu'une seule autre solution de nos
équations (en exigeant \marginal{des $p = \pm 1$ et} que $l_0^{p}$ et
$\mathrm{int}(g)\, l_1^{p}$ engendrent $\pi$, en plus des relations
fondamentales

\page{10}
plus haut, \struck{\ill{}} savoir $\tau$, correspondant \struck{\ill{}} à
\[
\begin{cases}
p = -1, \qquad g = l_0^{\alpha}\, l_1^{\beta} \quad (\alpha + \beta = 1), \quad \alpha, \beta \in \mathbb{Z}) \\
q = -1, \quad r = 1 \qquad \uncertain{au moins}, \quad \beta = -1
\end{cases}
\]

\note{« savoir $\tau$ » est porté au-dessus d'un passage biffé ; la fin de la
seconde ligne du système se lit mal}

\textbf{NB} Si on remplace dans les calculs qui précèdent $L_0$ par $L_1$, alors
$g$ est remplacé par $\rho(g)$, $q$ et $r$ restent inchangés. On trouve
ainsi, en plus de la \uncertain{fonction} $\chi$, \struck{\ill{}} \ill{}
fonctions \ill{} $\beta$, sur le groupe
\[
\mathrm{Autext}_{\mathrm{lac}}(\hat\pi_{0,3}, \mathfrak{S}_3) \overset{\text{déf}}{=} \Pi,
\]
\struck{\ill{}} $\gamma$ à valeurs dans $\hat{\mathbb{Z}}$, et une autre à valeurs
dans $\pi/\hat{\mathbb{Z}}$, en faisant opérer $\hat{\mathbb{Z}}$ sur $\pi$ par une
opération $\varphi$
\[
\varphi(\alpha).g = l_0^{\alpha}\, g\, l_1^{-\alpha}
\]

\note{« lac » est écrit en indice sous « Aut ext » ; au-dessus du premier
argument, un $\pi$ et un signe « $\|$ » qui le relient au $\hat\pi_{0,3}$}

L'application
\[
(\gamma, \chi) : \Pi \longrightarrow \pi/\hat{\mathbb{Z}} \times \hat{\mathbb{Z}}^{*},
\qquad u \longmapsto (\gamma(u), \chi(u))
\]
est injective, \struck{\ill{}} $\beta$ étant défini par $\gamma$ \ill{} $\gamma$
prend ses valeurs dans les classes des $g \in \pi$ (mod $\varphi(\hat{\mathbb{Z}})$)
satisfaisant \uncertain{une} relation

\marginal{NB $\gamma : E_0 \to \pi$ \ill{} cocycle sur \uncertain{$\Pi$} \ill{}
$(\pi, L_0)$ \ill{} \uncertain{1-cocycle} \ill{}}

\note{la note marginale est écrite en oblique dans la marge gauche, en
regard des trois dernières lignes ; elle ne se lit que par fragments}

\page{11}

\[
\sigma_\infty(g^{-1}) = l_0^{\beta}\, g\, l_1^{\beta}.
\]

Ce qui détermine $\beta$ de façon unique en fonction de $g$.

\struck{Appliquant la relation fondamentale $ug = \mathrm{int}(g)\, \uncertain{\rho}\, u$
à $l_0$, on trouve, puisque $u(l_0) = l_0$, $\rho(l_0) = l_1$,}
\[
\struck{u(l_1) = \mathrm{int}(g)(l_1)}
\]
\struck{et ceci détermine \ill{}} \struck{\ill{}} \struck{une \uncertain{multiplication} à droite
par une \uncertain{puissance} de $l_1$, près, le donnée de $u$}

\note{ce passage (de « Appliquant » à « le donnée de $u$ ») est barré de deux
grandes croix ; « le donnée » est tel quel}

On a de plus jolis générateurs de \add{la} \uncertain{sous-}extension
dans $\mathrm{Autext}(\pi)$, de $\mathfrak{S}_3$ par $\pi$ que les six générateurs
$l_0, l_1, l_\infty, \sigma_0, \sigma_1, \sigma_\infty$ (donnant $\rho$ par
$\sigma_0\sigma_1 = l_1\rho$ \uncertain{par ex.} i.e.
$\rho = \sigma_0\sigma_1 l_1$ p. ex.), en introduisant les \ill{} de
\uncertain{ramification} dans les pts $0, 1, \infty$, qui \ill{}
\uncertain{conjugués} \ill{} : $\hat{\mathbb{Z}}$,

\note{« sous- » et « dans $\mathrm{Autext}(\pi)$, » sont ajoutés au-dessus de
la ligne ; « $\mathrm{Autext}$ » se lit « Aut ext »}

On définira donc $\varepsilon_0, \varepsilon_1, \varepsilon_\infty$ par les relations
\[
\varepsilon_0^2 = l_0, \quad \varepsilon_1^2 = l_1, \quad \varepsilon_\infty^2 = l_\infty
\]
et plus précisément, $\varepsilon_i$ sera l'unique \struck{\ill{}} \ill{}
dans le \ill{} $\pi$ de $\sigma_i$, et qui fixe $l_i$ et est \ill{}

\page{12}

comme \uncertain{$\mathrm{int}(l_0)\sigma_\infty$} fixe $l_\infty$ (mais aussi
peut-être \uncertain{par} \ill{} avec $l_\infty$), on doit \ill{} multiplier
à droite par un $\mathrm{int}(l_\infty^{\alpha})$ $(\alpha \in \hat{\mathbb{Z}})$.

\[
\struck{\varepsilon_\infty(l_\infty) = l_\infty, \quad \varepsilon_\infty(l_0) = }
\]
\[
\varepsilon_\infty = \bigl(\mathrm{int}(l_0)\sigma_\infty\bigr)^{\mathrm{int}(l_\infty^{\alpha})}
\quad \text{i.e.} \quad
\varepsilon_\infty(l_0) = l_0 l_1 l_0^{-1}, \quad
\varepsilon_\infty(l_1) = l_0, \quad
\varepsilon_\infty(l_\infty) = l_\infty
\]

\note{« $\mathrm{int}(l_\infty^{\alpha})$ » est écrit en exposant au-dessus de la
parenthèse ; la lecture comme exposant est douteuse. Une accolade relie
la ligne biffée qui précède aux lignes suivantes}

\struck{et de même}
\[
\struck{\varepsilon_0 = \mathrm{int}(l_1)\sigma_0, \qquad
\varepsilon_1 = \mathrm{int}(l_\infty)\sigma_1}
\]

On \struck{aura} \ill{} \uncertain{Écrivant} que \uncertain{$g$} \ill{} au
$\mathrm{int}(g)$ pour $g \in \pi$, i.e. identifiant $\pi$ à un sous-groupe de
$\mathrm{Aut}(\pi)$, on trouve \struck{comme $\varepsilon_\infty = l_0\sigma_\infty$}
\[
\varepsilon_\infty = l_0\, \sigma_\infty\, l_\infty^{\alpha}.
\]
\marginal{doit écrire $\varepsilon_0^2 = l_\infty$ \ill{}}

\note{la note « doit écrire $\varepsilon_0^2 = l_\infty$ », interlinéaire,
remplace un mot biffé ; l'indice $0$ se lit mal, et la relation cherchée
plus bas est $\varepsilon_\infty^2 = l_\infty$}

\[
\struck{\varepsilon_\infty^2 = l_0\sigma_\infty l_0 \sigma_\infty
= l_0 \sigma_\infty l_0 \sigma_\infty^{-1} = l_0\, \sigma_\infty(l_0)}
\]
donc
\[
\begin{aligned}
\varepsilon_\infty^2 &= l_0\, \underbrace{\sigma_\infty\, l_\infty^{\alpha}\, l_0\, \sigma_\infty}\, l_\infty^{\alpha} \\
&= l_0\, \sigma_\infty(l_\infty^{\alpha} l_0)\, l_\infty^{\alpha}
= l_0\, (l_0^{-1} l_1^{-1})^{\alpha}\, l_1\, (l_1^{-1} l_0^{-1})^{\alpha} \\
&= \struck{l_0\, \sigma_\infty(l_\infty)\, \sigma_\infty(l_0)^{\alpha-1}\, \sigma_\infty(l_0)\, l_\infty^{\alpha}}
\end{aligned}
\]
\note{sous la ligne biffée, deux ébauches également biffées :
$l_0 l_\infty$ et $l_1^{-1} l_0^{-1}$}

et la relation $\varepsilon_\infty^2 = l_\infty$ s'écrit \uncertain{ainsi}, en
multipliant par $l_0 l_1$

\[
\struck{l_0 l_\infty\, \sigma_\infty(l_\infty)^{\alpha-1}\, \sigma_\infty(l_0)\, l_\infty^{\alpha-1} = 1}
\]
\[
\struck{\text{i.e.}\quad l_0^{-1} l_1\; l_0^{1-\alpha}\, l_1^{2-\alpha}\, l_\infty^{\alpha} = 1}
\]

\note{ces deux lignes sont barrées de traits obliques ; sous la première,
les accolades donnent $(l_0^{-1} l_1^{-1})^{\alpha-1}$ et $l_1$ pour les
deux facteurs soulignés}

\[
l_0 l_1 l_0\, (l_0^{-1} l_1^{-1})^{\alpha}\, l_1\, (l_1^{-1} l_0^{-1})^{\alpha} = 1
\]
\marginal{prenons $\alpha = 1$ ?}

\[
\struck{l_0 l_1 l_0 l_0^{-1} l_1^{-1} l_0^{-1} l_1^{-1} l_1 l_1^{-1} l_0^{-1} l_1 l_0^{-1} = 1 \quad ? \quad \text{Non !}}
\]
\[
\struck{l_0 l_1 l_0 l_0^{-1} l_1^{-1} l_1 l_1^{-1} l_0^{-1} = 1} \quad \text{ok}
\]

\note{les deux dernières lignes sont barrées lettre par lettre de traits
obliques (la première, en plus, d'un trait horizontal) ; le « ok » qui suit
la seconde n'est pas barré, et « donc » enchaîne}

donc
\[
\begin{cases}
\varepsilon_\infty = l_0\, \sigma_\infty\, l_\infty & \bigl(\text{ou encore } \mathrm{int}(l_0)\,\sigma_\infty\, \mathrm{int}(l_\infty)\bigr) \\
\varepsilon_0 = l_1\, \sigma_0\, l_0 \\
\varepsilon_1 = l_\infty\, \sigma_1\, l_1
\end{cases}
\]

\page{13}

donc
\[
\varepsilon_0^2 = l_0, \qquad \varepsilon_1^2 = l_1, \qquad \varepsilon_\infty^2 = l_\infty
\]

Ceux qui redonnent les $l_i$, et aussi les $\sigma_i$, \uncertain{par} en
termes des $\varepsilon_i$
\[
\sigma_0 = \varepsilon_0\, \varepsilon_1^2, \qquad
\sigma_1 = \varepsilon_1\, \varepsilon_\infty^2, \qquad
\sigma_\infty = \varepsilon_\infty\, \varepsilon_0^2
\]

D'où les relations
\[
\begin{cases}
\varepsilon_0^2\, \varepsilon_1^2\, \varepsilon_\infty^2 = 1 \\
(\varepsilon_0\, \varepsilon_1^2)^2 = 1 \\
(\varepsilon_1\, \varepsilon_\infty^2)^2 = 1 \\
(\varepsilon_\infty\, \varepsilon_0^2)^2 = 1
\end{cases}
\]

\uncertain{Pour obtenir} $\rho$ \struck{\ill{}} par
\[
\sigma_0 \sigma_1 = l_1^{-1} \rho \qquad \text{donc} \qquad
\rho = l_1 \sigma_0 \sigma_1
= \varepsilon_1^2\, \varepsilon_0\, \varepsilon_1^2\, \varepsilon_1\, \varepsilon_\infty^2
= \varepsilon_1^2\, \varepsilon_0\, \varepsilon_1\, \underbrace{\varepsilon_1^2\, \varepsilon_\infty^2}_{\varepsilon_0^{-2}}
\]
\[
\struck{\rho = \sigma_0\sigma_1 = \sigma_1\sigma_\infty = \sigma_\infty\sigma_0}
\qquad\qquad
\struck{\varepsilon_0\, \varepsilon_1^3\, \varepsilon_{\ill{}}}
\]

\note{la ligne biffée est suivie, sous son premier membre, d'un signe
« $\|$ » et de l'expression $\varepsilon_0\varepsilon_1^3\varepsilon_{\ldots}$,
biffée elle aussi}

\[
\begin{aligned}
\rho &= \varepsilon_1^2\, \varepsilon_0\, \varepsilon_1\, \varepsilon_0^{-2} \\
&= \varepsilon_\infty^2\, \varepsilon_1\, \varepsilon_\infty\, \varepsilon_1^{-2} \\
&= \varepsilon_0^2\, \varepsilon_\infty\, \varepsilon_0\, \varepsilon_\infty^{-2}.
\end{aligned}
\]
\marginal{ces relations se \uncertain{déduisent} \ill{} par \uncertain{conjugaison}
\uncertain{des autres}}

\note{la remarque marginale est séparée des trois égalités par un trait
vertical}

\uncertain{Multipliant} \struck{dans} ces relations dans l'ordre inverse et
écrivant $\rho^3 = 1$, on trouve
\[
(\varepsilon_0\, \varepsilon_1\, \varepsilon_\infty)^2 = 1
\]
ce qui est inquiétant, car il n'y a pas d'élément

\page{14}

d'ordre fini dans le groupe \uncertain{discret} \ill{} qu'il
\marginal{\struck{\ill{}} fixé par $\rho$}
faudrait donc avoir
\[
\varepsilon_0\, \varepsilon_1\, \varepsilon_\infty = 1
\]
i.e. en termes des $l_i$, $\sigma_i$
\[
l_1 \underbrace{\sigma_0\, l_0}\, l_\infty\, \underbrace{\sigma_1\, l_1}\, l_0\, \underbrace{\sigma_\infty\, l_\infty} = 1
\qquad \text{i.e.} \qquad l_\infty^{-1}\, l_0\, l_1^{-1} = 1 \quad \text{ok}
\]

\note{sous les trois accolades : $\struck{l_1 \sigma_\infty}$ ou
$\struck{l_1^{-1} l_\infty^{-1} \sigma_0}$ (deux essais, biffés), puis
$l_\infty^{-1} l_0^{-1}$ et $l_0^{-1} l_1^{-1}$. Le « ok » se lit sous la ligne,
à droite}

On a bien
\[
\boxed{\varepsilon_0\, \varepsilon_1\, \varepsilon_\infty = 1}
\]

On peut en déduire de plus jolies expressions pour les $\sigma_i$
\[
\sigma_0 = \varepsilon_\infty^{-1} \varepsilon_1, \qquad
\sigma_1 = \varepsilon_0^{-1} \varepsilon_\infty, \qquad
\sigma_\infty = \varepsilon_1^{-1} \varepsilon_0
\]
donc aussi la relation inattendue
\[
\sigma_\infty\, \sigma_1\, \sigma_0 = 1
\]

\note{tout ce qui précède, depuis le haut de la page jusqu'à « la relation
inattendue $\sigma_\infty\sigma_1\sigma_0 = 1$ », est barré de deux longs
traits obliques}

Sans doute $\varepsilon_0 \varepsilon_1 \varepsilon_\infty = \uncertain{\sigma_1}$
\[
\text{i.e.} \qquad
\varepsilon_0 = \varepsilon_1\, \varepsilon_\infty\, \varepsilon_1^{-1}, \qquad
\varepsilon_1 = \varepsilon_\infty\, \varepsilon_0\, \varepsilon_\infty^{-1}, \qquad
\varepsilon_\infty = \varepsilon_0\, \varepsilon_1\, \varepsilon_0^{-1}
\]

\note{le second membre de « Sans doute » est un signe court, lu $\sigma_1$ sans
conviction ; les trois relations qui suivent sont celles qu'on obtient en
supposant que les $\varepsilon_i$ se conjuguent circulairement}

\section{Critique des notations}

\note{ce titre est de l'édition ; il reprend les premiers mots du
paragraphe}

Critique des notations \uncertain{et} de l'approche. Les éléments $l_i$,
$\sigma_i$, $\varepsilon_i$, $\rho$ … sont dans le groupe discret, qu'il
faudrait distinguer du groupe profini — de plus il faudrait pour le groupe
discret \struck{utiliser} \ill{} pour le

\page{15}

profini, dans une certaine mesure !) utiliser le langage géométrique du
groupoïde fondamental de $\mathbb{P}^1(\mathbb{C}) - \{0, 1, \infty\}$.
\struck{On pose} On notera donc $\Pi = \Pi_{0,3}$ le groupe discret
\[
\Pi = \pi_1\bigl(\overbrace{\mathbb{P}^1(\mathbb{C}) - \{0, 1, \infty\}}^{U = U_{0,3}},\, P\bigr)
\]
où $\mathbb{P}^1(\mathbb{C})$ est muni de la structure \marginal{de sphère}
riemannienne \marginal{orientée} unique pour laquelle la rotation $\rho$ d'angle
$+\frac{2\pi}{3}$ autour de l'axe $P, \overline{P}$ \struck{\ill{}} induit sur
$(0, 1, \infty)$ la permutation circulaire $0 \to 1 \to \infty \to 0$ (ceci
correspond à l'orientation habituelle de la sphère), et où les
$\sigma_0, \sigma_1, \sigma_\infty$ \struck{\ill{}} sont les automorphismes
\ill{} de $\mathbb{P}^1_{\mathbb{C}}$ qui \struck{\ill{}} \marginal{\struck{fixent}}
induisent sur $(0, 1, \infty)$ les transpositions fixant resp. $0, 1, \infty$,
i.e. les \marginal{symétries} \struck{rotations} autour des axes $(0, 2)$,
$(1, -1)$ et $(\infty, \frac{1}{2})$. Quant à la conjugaison complexe, c'est la
symétrie p.r. au plan $(0, 1, \infty)$. On considère $\mathfrak{S}_3$ comme
opérant sur $(0, 1, \infty)$ dans cet ordre, et
$\mathfrak{S}_3 \times \mathbb{Z}/2$ comme opérant sur la sphère, via les
$\sigma_0, \sigma_1, \sigma_\infty \in \mathfrak{S}_3$ et
$\tau \in \mathbb{Z}/2$. On considère l'extension \struck{\ill{}} de
$\mathfrak{S}_3 \times \mathbb{Z}/2$ par $\Pi$, notée $\mathcal{E}$
\marginal{$\mathcal{E}_{0,3}$ ou}, définie par
\[
\mathcal{E} = \pi_1\bigl(\underset{\substack{\| \\ U_{0,3}}}{U},\;
\underbrace{\mathfrak{S}_3 \times \mathbb{Z}/2}_{\Gamma};\; P\bigr)
\]

\note{« de sphère », « orientée », « symétries » et « $\mathcal{E}_{0,3}$ ou »
sont des ajouts interlinéaires ; un mot biffé en oblique dans la marge
gauche, en regard de « Quant à la conjugaison », ne se lit pas}

\marginal{sur la sphère on a
$\sigma_0(z) = \frac{-z}{1-z} = \frac{z}{z-1}$,
$\sigma_1(z) = \frac{1}{z}$,
$\sigma_\infty(z) = 1 - z$,
$\rho(z) = \frac{1}{1-z}$,
$\rho^2(z) = 1 - \frac{1}{z}$}

\note{cette note est écrite en oblique dans la marge gauche, en regard de
« $(1,-1)$ et $(\infty, \frac{1}{2})$ » et des lignes suivantes}

\page{16}
\note{p. 7 de l'auteur}

Le groupe fondamental mixte est défini à la Malgoire, \struck{correspondant}
les classes $(\gamma, \lambda)$ d'un $\gamma \in \Gamma$ et d'une classe de chemins
$l$/de $P$ dans $\gamma^{-1}(P)$, en définissant
\[
(\gamma, l)(\gamma', l') = \bigl(\gamma\gamma',\; \gamma'^{-1}(l) \circ l'\bigr)
\]
i.e. « composant les chemins comme des fonctions » (c'est la convention pour
faire opérer à gauche le groupoïde fondamental sur les fibres
\uncertain{d'un} \uncertain{rev. étale} $U'/U$, à une classe de chemins
$l$ de $x$ à $y$ étant associé $l_* : U'_x \to U'_y$ …)

\note{la classe est notée $(\gamma, \lambda)$ puis $(\gamma, l)$ ; dans la
formule, « $\gamma\gamma'$ » est ajouté au-dessus de la première composante}

On choisit des éléments
\[
\dot\sigma_0, \dot\sigma_1, \dot\sigma_\infty \quad \text{sur} \quad
\sigma_0, \sigma_1, \sigma_\infty
\]
comme les couples d'un $\sigma_i$, et du chemin/des chemins de
$P$ à $\overline{P} = \sigma_i P$ qui se fait sur \struck{\ill{}} celui des
\uncertain{demi-grands} cercles de $P$ à $\overline{P}$, \struck{\ill{}} qui est
fixé par la symétrie $\tau\sigma_i = \sigma_i\tau$ (symétrie p.r. au plan
passant par $P, \overline{P}, a_i$ $(= 0, 1 \ldots \infty)$ et échangeant les
deux autres $a_j, a_k$), et \uncertain{un} \ill{} \uncertain{pour $1(-1)$}.
\struck{\ill{}} Désignant par $l_0, l_1, l_\infty$ les lacets \struck{\ill{}}
par $a_i$. Désignant par $l_0, l_1, l_\infty$ les lacets en sens direct autour
de $0, 1, \infty$ à partir de $P$, en allant vers $0, 1, \infty$ suivant des
segments de grands cercles, on trouve
\[
(1) \qquad l_\infty\, l_1\, l_0 = 1
\]
(et non $l_0 l_1 l_\infty = 1$ ! \uncertain{Voir} par l'écriture
$l_0^{-1} l_1^{-1} l_\infty^{-1} = 1$) et les relations

\note{le passage de « qui se fait sur » à « $a_j, a_k$ » est très corrigé ;
« celui des » est ajouté au-dessus d'un mot biffé. Les mots « pour $1(-1)$ »
sont entre parenthèses en bout de ligne, et la phrase « Désignant par
$l_0, l_1, l_\infty$ les lacets … par $a_i$ » est laissée inachevée avant
d'être reprise}

\page{17}

\[
(2) \qquad
\begin{cases}
\sigma_0^2 = \sigma_1^2 = \sigma_\infty^2 = 1 \\
\sigma_0\sigma_1 = \rho\, l_0, \qquad \sigma_1\sigma_\infty = \rho\, l_1, \qquad \sigma_\infty\sigma_0 = \rho\, l_\infty
\end{cases}
\]

\note{les seconds membres sont récrits : sous chacun on devine une première
écriture biffée, $l_0 \rho$, $l_1\rho$ \uncertain{} et $l_\infty\rho$, et
$\rho\, l_0$, $\rho\, l_1$, $\rho\, l_\infty$ sont portés au-dessus ou
au-dessous}

où $\rho$ \uncertain{au-dessus} de $\rho$, est défini par le chemin
direct \uncertain{de} $P$ à $P$. On a donc
\[
\begin{cases}
\rho\sigma_0\rho^{-1} = \sigma_1 \\
\rho\sigma_1\rho^{-1} = \sigma_\infty \\
\rho\sigma_\infty\rho^{-1} = \sigma_0
\end{cases}
\qquad \text{i.e.} \qquad
\begin{cases}
\rho\sigma_0 = \sigma_1\rho \\
\rho\sigma_1 = \sigma_\infty\rho \\
\rho\sigma_\infty = \sigma_0\rho
\end{cases}
\]

\marginal{NB Comme $\sigma_1 = \rho\sigma_0\rho^{-1}$, ces relations
s'écrivent aussi \struck{\ill{}}
$l_0 = \rho^{-1}\sigma_0\rho\sigma_0\rho^{-1}$, \struck{\ill{}}
$= \rho^{-1}[\sigma_0, \rho]$ \ill{} —
$\sigma_0(l_0) = l_0^{-1} l_1^{-1}$, $\sigma_0(l_1) = l_\infty$,
$\sigma_0(l_\infty) = l_1$}

\note{la note marginale occupe la marge droite, séparée du texte par un trait
vertical ; ses premières lignes sont surchargées et en partie hachurées}

On ne voit pas apparaître d'expression de $\rho$ en fonction des
$\sigma_0, \sigma_1, \sigma_\infty$ seuls — \struck{on voit en \ill{} (utilisant
le fait que $l_0, l_1, l_\infty$ \ill{} $\sigma$ \ill{} avec les relations
fondamentales (1) sont une présentation de $\Pi$, que $\rho$ doit s'exprimer
en fonction des $l_i$, $\rho$ \uncertain{puissance} \ill{}
$(\sigma_0\sigma_1)^3$ (p. ex.)}
\[
\struck{\rho = l_0 l_1 l_\infty}
\]
\struck{\ill{} $\sigma_i^2 = 1$}, il faudrait que $\rho$ soit produit de termes
$\sigma_i$ \uncertain{pris} deux consécutifs et \uncertain{$\sigma$}
\struck{\uncertain{abstraitement} \ill{} $\sigma_0\sigma_1\sigma_0 =$},

\note{le passage de « on voit » à « il faudrait » est barré d'une série de
traits obliques ; « deux consécutifs » est souligné}

\[
\begin{cases}
\dot\sigma_0^2 = \dot\sigma_1^2 = \dot\sigma_\infty^2 = 1 \\
\dot\sigma_0\dot\sigma_1 = \dot\sigma_1\dot\sigma_\infty = \dot\sigma_\infty\dot\sigma_0 = \dot\rho \\
\dot\rho^3 = 1
\end{cases}
\]
\marginal{\uncertain{\ill{} en \ill{}}}
est une présentation de $\mathfrak{S}_3$), que les symétries des relations
(1)(2) et \uncertain{autres}
\[
\rho^3 = 1
\]
\uncertain{donnent} une présentation de l'image inverse de $\mathfrak{S}_3$ dans
$\mathcal{E}$ (ou dans $\widehat{\mathcal{E}}$, au pt de vue profini).

\marginal{C'est faux (+) comme on \uncertain{voit} \ill{} \ill{} relations \ill{}
$\rho = 1$, \ill{} $\sigma_0\rho(\uncertain{?}) = \ill{}$}

\note{la note marginale est écrite en oblique dans l'angle inférieur gauche ;
elle ne se lit que par fragments. Le signe (+) renvoie peut-être à la
dernière phrase}

\page{18}

\[
(8) \qquad
\begin{cases}
\varepsilon_0 = l_0^{\ill{}}\, \sigma_0\, l_1^{\ill{}} \\
\varepsilon_1 = l_1^{\ill{}}\, \sigma_1\, l_\infty^{\ill{}} \\
\varepsilon_\infty = l_\infty\, \sigma_\infty\, l_0^{\ill{}}
\end{cases}
\qquad\qquad
\text{NB} \quad \rho\,\varepsilon_i\,\rho^{-1} = \varepsilon_{\rho(i)}
\quad \text{i.e.} \quad
\begin{cases}
\rho\varepsilon_0\rho^{-1} = \varepsilon_1 \\
\rho\varepsilon_1\rho^{-1} = \varepsilon_\infty \\
\rho\varepsilon_\infty\rho^{-1} = \varepsilon_0
\end{cases}
\]

\note{dans (8), les lettres et les indices sont repassés à l'encre épaisse et
chaque $l$ porte en exposant un petit signe surchargé qui ne se lit pas}

On a bien
\[
\varepsilon_0^2 = l_0\, \underbrace{\sigma_0\, l_1\, l_0\, \sigma_0}\, l_1
= l_0\, l_\infty\, l_\infty^{-1}\, l_1^{-1}\, l_1 = l_0 \qquad \text{ok}
\]
\[
\sigma_0(l_1 l_0) = \sigma_0(l_1)\, \sigma_0(l_0) = l_\infty\, (l_1 l_\infty)^{-1} = l_1^{-1}
\]

\note{le second facteur $l_0$ et le dernier $l_1$ de la première ligne sont
surchargés}

\textbf{NB} Les $\varepsilon_i$ ici sont les $\varepsilon_i^{-1}$ de tantôt !
\struck{qu'il faudrait \ill{} noter $\varepsilon'_i$ …})

Ceci dit, \uncertain{on a une} présentation de
$\mathcal{E}^{+}$ par générateurs et relations,
$l_0, l_1, l_\infty, \varepsilon_0, \varepsilon_1, \varepsilon_\infty, \rho$
\struck{de $\mathcal{E}^{+}$}, les relations
\[
\begin{cases}
l_\infty\, l_1\, l_0 = 1 \\
\varepsilon_0^2 = l_0, \quad \varepsilon_1^2 = l_1, \quad \varepsilon_\infty^2 = l_\infty \\
\varepsilon_0\varepsilon_1 = \rho(\ldots), \quad \varepsilon_1\varepsilon_\infty = \ldots, \quad \varepsilon_\infty\varepsilon_0 = \ldots, \quad \rho^3 = 1
\end{cases}
\]
où il faut compléter l'expression de $\varepsilon_0\varepsilon_1$ etc comme
produit de $\rho$ par des expressions en les $l_i$.
\struck{$\varepsilon_0\varepsilon_1 = l_0\sigma l_1$}

\[
\varepsilon_0\varepsilon_1 = l_0\sigma_0 l_1\, l_1\sigma_1 l_\infty
= l_0\, \underbrace{\sigma_0\, l_1^2\, \sigma_0^{-1}}_{l_\infty^2}\,
\underbrace{\sigma_0\sigma_1}_{\rho\, l_0}\, l_\infty
= l_0\, l_\infty^2\, \rho\, l_0\, l_\infty
\]

donc finalement
\[
(9) \qquad
\begin{cases}
\varepsilon_0\varepsilon_1 = l_0\, l_\infty^2\, \rho\, l_0\, l_\infty \\
\varepsilon_1\varepsilon_\infty = l_1\, l_0^2\, \rho\, l_1\, l_0 \\
\varepsilon_\infty\varepsilon_0 = l_\infty\, l_1^2\, \rho\, l_\infty\, l_1
\end{cases}
\]

\note{dans (9), les $\rho$ sont écrits sur d'autres lettres, et le premier
membre de la première ligne est repassé}

Éliminant les $l_i$ comme $\varepsilon_i^2$, on trouve des gén.
$\varepsilon_0, \varepsilon_1, \varepsilon_\infty, \rho$, une relation de
définition

\page{19}
\note{p. 8 de l'auteur}

Il revient au même de dire que $\sigma_0, \sigma_1, \sigma_\infty, \rho$ sont des
générateurs, avec relations de définition
\[
\begin{aligned}
&(4) \qquad
\overbrace{(\rho^{-1}\sigma_\infty\sigma_0 \struck{\ill{}})}^{l_\infty}\,
\overbrace{(\rho^{-1}\sigma_1\sigma_\infty \struck{\ill{}})}^{l_1}\,
\overbrace{(\rho^{-1}\sigma_0\sigma_1 \struck{\ill{}})}^{l_0} = 1 \\
&(5) \qquad \sigma_0^2 = \sigma_1^2 = \sigma_\infty^2 = 1 \\
&(6) \qquad \rho^3 = 1.
\end{aligned}
\]

\marginal{NB $\sigma_0(l_0) = l_\infty^{-1} l_1^{-1}$, $\sigma_0(l_1) = l_\infty$,
$\sigma_0(l_\infty) = l_1$ \struck{\ill{}}, et itou pour $\sigma_1, \sigma_\infty$}

\textbf{NB} Si on y fait \add{de plus} $\sigma_0, \sigma_1, \sigma_\infty = 1$,
on trouve comme relation de définition $\rho^3 = 1$, \struck{\ill{}} sans plus,
ce qui prouve \uncertain{aussi} que $\rho$ \uncertain{n'est pas}
\ill{} \uncertain{p.r.} aux $\sigma_i$ …

\note{l'ajout « de plus » est écrit au-dessus de « $\sigma_0, \sigma_1,
\sigma_\infty$ » ; « $= 1$ » est en bout de formule}

Tout ceci est compatible avec les relations des \uncertain{listes},
\struck{\ill{} dans ces dernières, on \ill{} par \ill{}, et qu'on interprète
$l'_i = l_i^{-1}$, $\sigma'_i = \sigma_{\rho(i)}$} \struck{i.e. $\sigma_i$}
\struck{\ill{}} si on y remplace les $l_i$ par les $l_i^{-1}$.

Ceci dit, comme $\sigma_i = \sigma^{-1}_{\uncertain{\rho(i)}}$
\[
(7) \qquad
\begin{cases}
\mathrm{int}(\sigma_i) : \text{échange } l_j \text{ et } l_k \quad (i, j, k \text{ distincts}) \\
\mathrm{int}(\rho)\, l_i = l_{\rho(i)}
\end{cases}
\]
\struck{(et \ill{} $l_{\rho(i)}$, \ill{} déduit par (1))}

\note{« $(i, j, k$ distincts$)$ » est porté au-dessus d'un passage biffé ; le
« comme $\sigma_i = \ldots$ » qui introduit (7) est mal lu}

On voudrait maintenant expliciter les $\varepsilon_i$
\[
\text{avec} \qquad \varepsilon_i^2 = l_i \quad (i = 0, 1, \infty),
\]
$\varepsilon_i$ étant \uncertain{dans} un générateur des automorphismes (ou
normalisateurs, c'est à dire) des $L_i$ dans $\mathcal{E}^{+}$
(\uncertain{ou} le normalisateur dans $\mathcal{E}$). Les calculs faits plus
haut, avec les relations \uncertain{inexactes}, donnant
\[
\struck{(8) \quad
\begin{cases}
\varepsilon_0 = l_0^{\ill{}}\, \sigma_1\, l_{\ill{}}^{\ill{}} \\
\varepsilon_1 = l_\infty^{-1}\, \sigma_\infty\, l_1^{-1} \\
\varepsilon_\infty = l_0^{-1}\, \sigma_0\, l_\infty
\end{cases}}
\]

\note{ce premier (8) est barré ; la page 18 le récrit}

\marginal{Il faut \uncertain{ajouter} \uncertain{la} relation (4) \ill{}
$\rho\sigma_0\rho^{-1} = \sigma_1$ \ill{} relations (5) \ill{} Noter
\ill{} (7) \ill{}
$l_0 = \ill{}$, $l_1 = \ill{}$, $l_\infty = l_0^{\ill{}} l_1$,
$\mathrm{int}(\sigma_0) l_1 = l_\infty^{\ill{}}$, $\mathrm{int}(\sigma_0) l_\infty = l_1^{\ill{}}$}

\note{la marge gauche de la page est couverte, en oblique, de calculs serrés
et en partie hachurés, dont seuls ces fragments se lisent}

\page{20}

Donc aussi les relations
\[
(13) \qquad
\begin{cases}
\tilde\sigma_0^2 = \tilde\sigma_1^2 = \tilde\sigma_\infty^2 = 1 \\
\tilde\sigma_0\tilde\sigma_1 = \tilde\sigma_1\tilde\sigma_\infty = \tilde\sigma_\infty\tilde\sigma_0 = \rho\, \struck{(\ill{})} \\
(\rho^3 = 1 \text{ est conséquence des précédentes})
\end{cases}
\]
et l'image inverse $\widetilde{\Sigma}$ de $\varphi(\mathfrak{S}_3) = \Gamma$
\uncertain{(ss-groupe diédral 2)} est produit ½ direct de $\mathfrak{S}_3$ et de
$\Pi$. Il faut expliciter les opérations des $\tilde\sigma_i$ sur $\Pi$ pour
avoir la structure complète \uncertain{du} groupe, on trouve pour
$\tilde\sigma_i$ ci-dessus $\varphi(\tilde\sigma_i)$ i.e. la symétrie p.r. au
plan $(P, \overline{P}, a_i)$, qui transforme les lacets
$l_0^{-1}$, \struck{\ill{}} $l_1$ \uncertain{en} $l_\infty^{-1}$
\struck{\ill{}} :
\[
(14) \qquad
\begin{cases}
\tilde\sigma_0\, l_0\, \tilde\sigma_0^{-1} = l_0^{-1} \qquad
\tilde\sigma_0\, l_1\, \tilde\sigma_0^{-1} = l_\infty^{-1} \qquad
\tilde\sigma_0\, l_\infty\, \tilde\sigma_0^{-1} = l_1^{-1} \\
\tilde\sigma_1\, l_1\, \tilde\sigma_1^{-1} = l_1^{-1} \qquad \ldots \\
\ldots
\end{cases}
\]

\note{la phrase « qui transforme les lacets … » est surchargée et se lit mal
avant les deux points ; (14) la remplace}

Les relations (13) (14), plus la relation fondamentale
$l_\infty l_1 l_0 = 1$ (1), forment une présentation de $\widetilde{\Sigma}$,
dont on peut \uncertain{(si on veut)} éliminer $\rho$,
\uncertain{se contentant} d'écrire la relation qui le définit)

On peut maintenant \uncertain{revenir} à \uncertain{nos} trois
\ill{}, $\sigma$ \uncertain{par}
\[
(15) \qquad
\begin{cases}
\tau_0 = \sigma_0\tilde\sigma_0 = \tilde\sigma_0\sigma_0 \\
\tau_1 = \sigma_1\tilde\sigma_1 = \tilde\sigma_1\sigma_1 \\
\tau_\infty = \sigma_\infty\tilde\sigma_\infty = \tilde\sigma_\infty\sigma_\infty
\end{cases}
\qquad \text{donc} \qquad
(16) \qquad
\begin{cases}
\tau_0^2 = 1 \\
\tau_1^2 = 1 \\
\tau_\infty^2 = 1
\end{cases}
\]

\marginal{NB $\rho\tilde\sigma_i\rho^{-1} = \tilde\sigma_{\rho(i)}$ i.e.
$\rho\tilde\sigma_0\rho^{-1} = \tilde\sigma_1$ \ill{}
$\rho\tilde\sigma_\infty\rho^{-1} = \tilde\sigma_0$}

\marginal{ce \uncertain{sont elles} \ill{} \uncertain{transport de structure}
\ill{} des \ill{} $(\overline{P}, P)$}

\marginal{Critique \ill{} $\tilde\sigma_0$ \ill{} \uncertain{s-groupe} \ill{}
\ill{} il \uncertain{vaut mieux} écrire $\tilde\sigma_0$ \ill{}
\uncertain{plusieurs} \ill{}}

\note{la marge gauche porte, en oblique, trois notes de sa main ; seuls ces
fragments se lisent}
\end{document}
